% Einheit 2 von 4 – Rechte Winkel rückwärts (bank/orthogonalitaet/e2.jsonl, zusammenbau v0.1)
%% TODO Zeitmarke Einheit 2: Marken-Zeile nicht lesbar
%% TODO Zweigzeile Teil 1: was hier gelernt wird (Satz in Schülersprache aus dem Einheitstitel) – Einheit 2
\einheitenkopf[e2]{Einheit 2 von 4 · Rechte Winkel rückwärts}
\zweigzeile{Abitur GK}
\setcounter{aufgabe}{13}

%% TODO Titel: Ich-kann-Satz fehlt in der Bank (Platzhalter: Kettenname) – Nr. 14
%% TODO Anweisung: ein Satz über der ersten Teilaufgabe fehlt in der Bank – Nr. 14
\begin{aufgabe}{Rechte Winkel rückwärts}
%% TODO \rechenplatz{n}: Schrittzahl des Grundfalls fehlt in der Bank (nur bei fester Schreibform, 2.3 b)
\begin{teile}
\teil Aufgabe: „Bestimme $a$ so, dass das Dreieck $OAB$ bei $B$ rechtwinklig ist.“ Was ist zu tun? \\ \kreuz{null zeigen – für jeden Wert des Parameters} \\ \kreuz{null setzen – ein Wert ist gesucht, Gleichung lösen} \\ \kreuz{null zeigen – ohne Parameter, einsetzen und ausrechnen}
\teil Gegeben sind der Ursprung $O$, $A(4 | 0 | a)$ und $B(2 | 2 | 3)$. Bestimme $a$ so, dass das Dreieck $OAB$ bei $B$ rechtwinklig ist \hfill (Abitur 2021 GK).
\teil Der Ursprung $O$, $P(4 | 0 | 12)$ und $Q_t(2 | t | 8)$ bilden ein Dreieck. Bestimme alle Werte von $t$, für die das Dreieck bei $Q_t$ einen rechten Winkel hat \hfill (Abitur 2024 GK).
\teil Der Punkt $Q$ liegt auf der Kante von $A(4 | 0 | 0)$ nach $D(4 | 0 | 10)$; außerdem sind $R(0 | 3 | 6)$ und $F(10 | 0 | 8)$ gegeben. Das Dreieck $FQR$ soll bei $Q$ einen rechten Winkel haben. Berechne die $z$-Koordinate von $Q$ \hfill (Abitur 2023 LK).
\teil Ein Quader hat die Eckpunkte $O(0 | 0 | 0)$, $A(5 | 0 | 0)$, $B(5 | 12 | 0)$, $C(0 | 12 | 0)$ und $E(5 | 0 | h)$ mit $h > 0$; $G$ liegt über $C$. Die Raumdiagonalen $AG$ und $CE$ stehen senkrecht zueinander. Bestimme $h$ und das Volumen des Quaders \hfill (Abitur 2026 LK).
\teil Gegeben sind $A(1 | 2 | 0)$ und $B(4 | 4 | 2)$. Für den Punkt $C$ gilt: Seine $z$-Koordinate ist $5$, die Geraden $AB$ und $AC$ verlaufen senkrecht zueinander, und $C$ liegt in der Ebene $2x + y = 0$. Bestimme die Koordinaten von $C$ \hfill (Abitur 2026 GK).
\teil Gegeben sind $A(0 | 0 | 0)$, $B(3 | 4 | 0)$ und $C(25 | 0 | 0)$. Der Punkt $T$ liegt auf der Strecke $AC$ so, dass das Dreieck $ABT$ bei $B$ rechtwinklig ist. Bestimme das Verhältnis $|AT| : |CT|$ \hfill (Abitur 2021 LK).
\teil Das Dreieck $ABC$ mit $A(0 | 0 | 0)$ und $B(1 | 4 | 1)$ ist gleichschenklig und rechtwinklig mit dem rechten Winkel bei $A$; die Kathete $AC$ liegt in der $xz$-Ebene. Bestimme die Koordinaten eines möglichen Punktes $C$ \hfill (Abitur 2023 LK).
\end{teile}
\end{aufgabe}

%% TODO Titel: Ich-kann-Satz fehlt in der Bank (Platzhalter: Kettenname) – Nr. 15
%% TODO Anweisung: ein Satz über der ersten Teilaufgabe fehlt in der Bank – Nr. 15
\begin{aufgabe}{Rechte Winkel rückwärts – Fehler finden, Begründen, Anwendung}
\begin{teile}
\teil Für $O$, $A(4 | 0 | a)$ und $B(2 | 2 | 3)$ soll $a$ so bestimmt werden, dass das Dreieck $OAB$ bei $B$ rechtwinklig ist. Nele rechnet: \rechnung{\overrightarrow{OA} \circ \overrightarrow{OB} &= 8 + 0 + 3a = 0 \\ a &= -\tfrac{8}{3}} Finde den Fehler.
\teil Begründe, warum die Forderung „rechter Winkel bei $B$“ bei einer unbekannten Koordinate genau eine Gleichung liefert.
\teil Ein Regal wird durch den Quader mit $O(0 | 0 | 0)$, $A(9 | 0 | 0)$, $B(9 | 12 | 0)$, $C(0 | 12 | 0)$ und $E(9 | 0 | h)$ beschrieben (1 LE = 1 dm). Zwei Stahlseile verlaufen entlang der Raumdiagonalen $AG$ und $CE$ ($G$ über $C$) und sollen sich senkrecht kreuzen. Wie hoch muss das Regal sein, und wie viel Volumen hat es?
\end{teile}
\end{aufgabe}
